% ============================================================
%  Block A -- the local properties
% ============================================================
\rsection{Block A --- at a point}

\begin{signpost}[What is being compared]
Fix $p$, a limit point of the domain of $f$. Three properties are in play:
$\lim_{x\to p}f(x)$ exists; $f$ is continuous at $p$; $f$ is differentiable
at $p$. They form a strict chain, and the two strictnesses cost one function
each.
\end{signpost}

\ritem{T1}{Theorem}
\emph{If $f$ is differentiable at $p$ then $f$ is continuous at $p$.}

\begin{proof}
For $x \ne p$,
\[ f(x)-f(p) \;=\; \frac{f(x)-f(p)}{x-p}\,\cdot\,(x-p) . \]
As $x \to p$ the first factor tends to $f'(p)$ and the second to $0$, so the
product tends to $0$ by the algebra of limits (4.4). Hence
$f(x) \to f(p)$.
\end{proof}

\noindent
This is Rudin's 5.2, and the proof is worth one remark: it uses the existence
of $f'(p)$ only to know the quotient has \emph{some} finite limit.\kn{So the
theorem is really ``if the difference quotient is bounded near $p$ then $f$
is continuous at $p$'' --- differentiability is more than is needed. That
observation is what makes the Lipschitz condition of Chapter 4's Supplement
E the natural weakening.}

\ritem{T2}{Theorem}
\emph{If $f$ is continuous at $p$ then $\lim_{x\to p}f(x)$ exists and equals
$f(p)$.}

\begin{proof}
Definition 4.5 says precisely this whenever $p$ is a limit point of the
domain.
\end{proof}

\begin{pitfall}[The one place the chain is not a chain]
At an \emph{isolated} point of the domain, $f$ is continuous automatically
(4.6) while $\lim_{x\to p} f(x)$ is not defined at all --- Rudin's 4.1
requires $p$ to be a limit point. So T2 needs the hypothesis, and the arrow
in the diagram is an arrow only over limit points.

This is the deleted-versus-undeleted distinction of Chapter 4's Supplement A
(S4.2). Rudin's limit is deleted, which is exactly why continuity is not
literally the statement ``the limit equals the value'' --- at isolated points
there is no limit to equal it.
\end{pitfall}

\noindent
\textbf{Both converses fail.} $\mathbf{E1}$, the removable spike, has a limit
at $0$ and is not continuous there; $\mathbf{E2}$, $|x|$, is continuous at
$0$ and not differentiable there.

\ritem{T3}{Theorem}
\emph{(Darboux.) If $f$ is differentiable on $[a,b]$ then $f'$ takes every
value between $f'(a)$ and $f'(b)$ --- even though $f'$ need not be
continuous.}

\begin{proof}
Suppose $f'(a) < \lambda < f'(b)$ and put $g(x) = f(x) - \lambda x$, so
$g'(a) < 0 < g'(b)$. Since $g'(a)<0$, $g(t)<g(a)$ for some $t$ just right of
$a$; since $g'(b)>0$, $g(s)<g(b)$ for some $s$ just left of $b$. So the
minimum of the continuous $g$ on the compact $[a,b]$ (4.16) is attained at an
interior point $x$, and there $g'(x)=0$ by 5.8 --- that is, $f'(x)=\lambda$.
\end{proof}

\ritem{T4}{Corollary}
\emph{A derivative has no simple discontinuities. In particular the signum
function is not the derivative of anything.}

\begin{proof}
A simple discontinuity is one at which both one-sided limits exist
(Definition 4.26). If they existed and differed, $f'$ would omit the values
strictly between them on either side of the point, contradicting T3; if they
existed and agreed but differed from $f'(x)$, the same argument applies to a
value between $f'(x)$ and the common one-sided limit.
\end{proof}

\begin{deeper}[What T3 and T4 buy, read against Chapter 4's taxonomy]
Chapter 4's appendix opened Rudin's two kinds of discontinuity into four
(S4.19): removable, jump, asymptotic, oscillatory. T4 says a derivative may
suffer only the last two. It cannot jump, and it cannot have a removable
discontinuity.

$\mathbf{E3}$, the function $x^2\sin(1/x)$, shows the permitted case really
occurs: it is differentiable everywhere, and its derivative oscillates
without limit at $0$. So ``differentiable'' and ``continuously
differentiable'' are genuinely different hypotheses --- which is why Rudin
states them separately from Chapter 5 onward, and why $\mathscr{C}^1$ needs
its own name.
\end{deeper}

\ritem{T5}{Theorem}
\emph{If $\lim_{x\to p}f(x)$ exists then $f$ is bounded on some punctured
neighbourhood of $p$.}

\begin{proof}
Take $\varepsilon = 1$ in Definition 4.1: there is $\delta>0$ with
$|f(x)-q|<1$, hence $|f(x)| < |q|+1$, for all $x$ in the punctured
$\delta$-neighbourhood.
\end{proof}

\noindent
Boundedness is the weakest thing in this section, and it is strictly weaker
than everything above it: $\mathbf{E4}$, $\sin(1/x)$, is bounded everywhere
and has no limit at $0$.\kn{The chain therefore extends one step downward:
differentiable $\Rightarrow$ continuous $\Rightarrow$ has a limit
$\Rightarrow$ locally bounded, with all three implications strict. Rudin
never states T5, but he uses it --- at 4.4, where the product of limits needs
one factor bounded, and again at 5.2.}
